\clearpage
\phantomsection% 
\addcontentsline{toc}{chapter}{KATA PENGANTAR}
%\thispagestyle{fancy}

\begin{justifying}
	\large \bfseries \centering \MakeUppercase{Kata Pengantar}\linebreak

	\normalsize \normalfont \justifying
	Puji syukur kepada Tuhan Yang Maha Esa atas berkat, rahmat, dan penyertaan-Nya sehingga penulis dapat menyelesaikan tugas akhir ini dengan baik. Dalam penyusunan tugas akhir ini, penulis telah banyak mendapatkan arahan, bantuan, serta dukungan dari berbagai pihak. Oleh karena itu, pada kesempatan ini penulis mengucapkan terima kasih yang sebesar-besarnya kepada: \par
	\begin{enumerate}
		\item Untuk Ibuku tercinta, Marni Simanjuntak, wanita yang kasih sayang dan doanya selalu menjadi kekuatan dalam setiap langkahku. Mungkin tidak ada kata yang cukup untuk menggambarkan rasa terima kasih atas segala pengorbanan, kesabaran, dan cinta yang telah Ibu berikan selama ini. Di setiap keberhasilan yang kuraih, selalu ada doa Ibu yang mengiringinya. Jika ada satu harapan yang ingin kusampaikan, semoga Ibu selalu diberikan kesehatan, kebahagiaan, dan umur yang panjang agar dapat menyaksikan setiap mimpi yang selama ini Ibu perjuangkan untukku menjadi kenyataan.
		\item Untuk keluargaku tercinta, Kak Liza, Kak Mena, Kak Keke, Kak Cilin, Kak Inat, dan Nathan, yang selalu menjadi rumah dalam setiap perjalanan hidupku. Terima kasih atas dukungan, perhatian, dan kepercayaan yang selalu kalian berikan, bahkan ketika aku meragukan diriku sendiri. Kalian telah mengajarkanku arti kebersamaan, ketulusan, dan kekuatan keluarga. Semoga setiap langkah yang kita tempuh membawa kebahagiaan, keberhasilan, dan kesempatan untuk saling membanggakan satu sama lain.
		\item Bapak I Wayan Wiprayoga Wisesa, S.Kom., M.Kom. selaku Dosen Pembimbing atas ide, waktu, tenaga, perhatian, dan masukan yang telah disumbangsihkan kepada penulis selama proses penelitian dan penulisan tugas akhir ini.
		\item Untuk seseorang yang begitu istimewa, yang senyumnya mampu membawa ketenangan dan yang kehadirannya telah memberi warna dalam perjalanan hidupku. Mungkin kamu tidak pernah menyadarinya, tetapi di balik banyak langkah yang kuambil, ada semangat yang lahir dari keberadaanmu. Jika suatu hari kata-kata ini sampai kepadamu, aku hanya berharap semoga setiap impian yang kamu perjuangkan dapat terwujud dan setiap kebahagiaan yang kamu cari dapat menemukan jalannya. Dan meskipun mungkin terdengar sedikit egois, aku berharap suatu hari nanti aku bisa menjadi bagian dari alasan di balik senyummu yang paling tulus.
		\item Untuk sahabat-sahabatku, Zakhi Alghifari dan Machzaul Harmansyah, yang telah menjadi bagian penting dalam perjalanan ini. Terima kasih atas setiap bantuan, diskusi, tawa, dan dukungan yang membuat berbagai tantangan terasa lebih ringan untuk dihadapi. Semoga persahabatan yang telah terjalin tidak hanya menjadi bagian dari masa lalu yang indah, tetapi juga terus tumbuh seiring dengan perjalanan dan impian yang kita kejar masing-masing.
	\end{enumerate} \par
	Akhir kata, penulis berharap semoga tugas akhir ini dapat memberikan manfaat bagi kita semua. Penulis menyadari bahwa tugas akhir ini tidak luput dari kekurangan dan kelemahan, sehingga penulis terbuka untuk menerima saran, kritik, dan masukan yang membangun.
	\vfill

	\begin{flushright}
		Lampung Selatan, Mei 2026 \par
		\vspace{5em}
		Lucky Immanuel Sitanggang \par
		NIM 122140179
	\end{flushright}

\end{justifying}
\clearpage





\begin{comment}
\clearpage
\pagestyle{fancy}
\fancyhf{}
\fancyhead[R]{\thepage}
\phantomsection% 
%\clearpage
%\phantomsection% 
\addcontentsline{toc}{chapter}{KATA PENGANTAR}
%\thispagestyle{fancy}

\begin{justifying}
	\large \bfseries \centering \MakeUppercase{Kata Pengantar}\linebreak
	
	\normalsize \normalfont \justifying
	Puji syukur kehadirat Tuhan Yang Maha Esa atas limpahan rahmat, karunia, serta petunjuk-Nya sehingga penyusunan tugas akhir ini telah terselesaikan dengan baik. Dalam penyusunan tugas akhir ini penulis telah banyak mendapatkan arahan, bantuan, serta dukungan dari berbagai pihak. Oleh karena itu pada kesempatan ini penulis mengucapan terima kasih kepada: \par
	\begin{enumerate}
		\item Bapak Prof. Dr. I. Nyoman Pugeg Aryantha selaku Rektor Institut Teknologi Sumatera.  
		\item Bapak Hadi Teguh Yudistira, S.T., Ph.D. selaku Dekan Fakultas Teknologi Industri.
		\item Bapak Andika Setiawan, S. Kom., M. Cs. selaku Ketua Program Studi Teknik Informatika.
		\item Bapak Ilham Firman Ashari, S. Kom., M.T. selaku Koordinator Tugas Akhir Program Studi Teknik Informatika.  
		\item Bapak Martin C. T. Manullang, Ph.D. selaku Dosen Pembimbing atas ide, waktu, tenaga, perhatian, dan masukan yang telah disumbangsihkan kepada penulis.
            \item Bapak Andika Setiawan, S. Kom., M. Cs dan Bapak Eko Dwi Nugroho, S.Kom., M.Cs. selaku Dosen Penguji atas saran dan masukan yang diberikan. 
		\item Kedua Orang Tua dan Adik yang selalu memberikan dukungan dan doa sehingga penulis dapat menyelesaikan tugas akhir ini. Kelembutan, dukungan, dan cinta yang kalian berikan selalu menjadi sumber inspirasi dan kekuatan.
		\item Teman-teman penulis yang membantu selama masa perkuliahan yang tidak bisa disebutkan satu persatu.
	\end{enumerate} \par
	Akhir kata penulis berharap semoga tugas akhir ini dapat memberikan manfaat bagi kita semua. Penulis menyadari bahwa tugas akhir ini tidak luput dari kekurangan dan kelemahan, dan penulis  terbuka untuk menerima saran, kritik, dan masukan yang membangun.
	\vfill
	
\end{justifying}
\clearpage
\end{comment}
